% Written by moonshotai/kimi-k3.1-pro (Open Math OS), 2026-11-13; unreviewed AI draft.
\documentclass[10pt]{article}
\usepackage[margin=.85in]{geometry}
\usepackage{amsmath,amssymb,amsthm}
\usepackage{xcolor}
\usepackage[colorlinks=true,linkcolor=black]{hyperref}
\newtheorem{lemma}{Lemma}
\newtheorem{proposition}{Proposition}
\newtheorem{theorem}{Theorem}
\theoremstyle{definition}\newtheorem*{definition}{Definition}
\theoremstyle{remark}\newtheorem*{remark}{Remark}
\newcommand{\F}{\mathfrak F}
\newcommand{\GAP}[1]{\textcolor{red}{\textbf{GAP: }#1}}
\title{Every internal node is worth $-p$:\\the accounting proof for Fuss--Catalan numbers}
\author{Mighty Math Magicians --- unreviewed AI draft}
\date{2026-11-13}
\begin{document}\maketitle
\paragraph{Status.} This is an AI-written proof draft, not a human-signed theorem.
It derives a factorial valuation by symmetry; see Remark~\ref{qualification}.
Verification is in \texttt{experiments/E004\_accounting\_proof/}, with
qualifications in \texttt{GAPS\_accounting.md}. Red gaps, if any, make the
remaining statement conditional; they are not hidden by the computations.

Fix a prime $p$, $k\ge1$, and $q=p^k$. A $p$-ary tree is a leaf or an ordered
$p$-tuple of such trees. It has $N(m)=(m-1)/(p-1)$ internal nodes if it has $m$
leaves. Write $m=\sum_jm_jp^j$ in base $p$, $s=s_p(m)=\sum_jm_j$, and
\[
 \nu(m)=\frac{m-s}{p-1}=\sum_{j\ge1}\left\lfloor\frac m{p^j}\right\rfloor,
 \qquad t=N(m)-\nu(m)=\frac{s-1}{p-1}.
\]
The Fuss--Catalan number $C_n^{(q)}=\binom{qn}{n}/((q-1)n+1)$ counts ordered
full $q$-ary trees with $n$ internal nodes. Set $m=(q-1)n+1$. Expanding every
$q$-ary node into a complete $p$-ary block of height $k$ converts it to a
$p$-ary tree with every leaf depth divisible by $k$.

\section{L1 and L2: labelled trees and cherries}
\begin{definition} In a \emph{leaf-labelled cyclic $k$-blocked $p$-ary tree},
every leaf has one of distinct labels $1,\ldots,m$ and lies at depth divisible
by $k$; at every internal node the $p$ children form a cyclic sequence. Its
\emph{rotation class} is the equivalence class of ordered trees obtained by
rotating any nonempty family of child cycles.\end{definition}

\begin{lemma}[L1]\label{L1} A labelled rotation class has exactly $p^{N(m)}$
ordered representatives. If $L(m)$ is the number of classes, then
$m!C_n^{(q)}=p^{N(m)}L(m)$.\end{lemma}
\begin{proof} For one internal node, rotating a labelled ordered tree cannot
fix it: a root rotation which returned the original child word would force
the same leaf label to occur more than once. If a composition of local rotations
fixes the root, recursively restricting it to each child shows it is fixed there;
then the root rotation is itself trivial. Hence the root action is trivial too.
Inductively the product of one $\mathbb Z_p$ at each internal node acts freely
on ordered labelled trees. Its orbits are precisely the labelled rotation
classes. There are $m!C_n^{(q)}$ ordered labelled trees in total.\end{proof}

For general cyclic trees, each leaf also has \emph{class} $c\in\mathbb Z/k$;
admissibility requires its depth to be $\equiv-c\pmod k$.
\begin{lemma}\label{cherry} Cyclicly permuting $p$ distinct same-class labels
fixes precisely those admissible cyclic trees in which those labels form
$p$ sibling leaves. Collapsing them on both sides gives a bijection with $p-1$
copies of admissible trees with one fewer $p$-ary node and one class-$(c+1)$
leaf.\end{lemma}
\begin{proof} Let $\varphi$ be the induced automorphism and $v$ the lowest common
ancestor of the selected leaves. If $\varphi$ acts trivially at $v$, it fixes
each child position, and therefore preserves the child subtree containing
any selected leaf in that position. At most one child subtree could then
contain more than one selected leaf; since $p\ge2$, repeating the argument in
that subtree would keep the successive ancestor lower than $v$, contradiction.
If instead $\varphi$ changes the root child word by nontrivially cycling the
selected labels, every selected leaf beneath $v$ moves, while all unselected
leaves of the entire tree are fixed. Since every child of $v$ has a leaf, no
child containing an offspring remains that was not selected. Consequently
the children of $v$ are exactly the selected leaves. Once $x_1$ is placed
in slot $1$, the remaining labels are forced to be $x_2,\ldots,x_p$ in one of
the $p-1$ starting rotations. All of those are fixed. Conversely, a
class-$(c+1)$ leaf expands to a sibling $p$-tuple of class $c$, restoring
the original tree.\end{proof}

\begin{lemma}[L2]\label{L2} Let $\mathcal M(r)$ be the set of all unrestricted
cyclic $p$-ary trees with $r$ labelled leaves. Then
$|\mathcal M(r)|\equiv(-1)^{N(r)}\pmod p$, and hence
$L(m)\equiv(-1)^{N(m)}\pmod p$.\end{lemma}
\begin{proof} If $r=1$ the set is a singleton; otherwise $r\equiv1\pmod{p-1}$
is necessary. Apply Lemma~\ref{cherry} to fixed labels $1,\ldots,p$. Because
all orbits other than fixed points have size $p$,
$|\mathcal M(r)|\equiv(p-1)|\mathcal M(r-p+1)|\pmod p$. Iterate exactly
$N(r)$ times. For blocked trees choose the next $p$ of the still present labels;
all class shifts remain zero at every collapse, giving the same recursion.\end{proof}

\section{L3: complete-tree forests}
\begin{definition} A \emph{digit-forest structure} on $[m]$ is a collection of
$m_j$ cyclic complete $p$-ary trees of height $j$, all disjoint, with leaf sets
partitioning $[m]$. Let $\F(m)$ denote the set of these structures. A tree of
height zero is a labelled leaf; the empty forest on $[0]$ is one object.\end{definition}

\begin{lemma}\label{stabilizer} $S_m$ is transitive on $\F(m)$, and its
stabilizer is $\prod_j(C_p^{N(p^j)}\wr S_{m_j})$, of order
$p^{\nu(m)}\prod_jm_j!$. Therefore
\begin{equation}\label{digit}
 m!=p^{\nu(m)}\prod_jm_j!\,|\F(m)|.
\end{equation}\end{lemma}
\begin{proof} Lift each cyclic component by choosing a root slot, keeping
the child subtrees cyclic. This gives $p$ root choices, while equal-height
components are still an unordered collection. A full leaf labelling assigns
$[m]$ bijectively to the $m$ leaf slots, so it gives $m!$ such objects. By the
argument of Lemma~\ref{L1}, independent local rotations act freely; an
automorphism of an ordered labelled component fixes its root slot and then
all lower nodes. Thus this lifted labelled class has size
$p^{\nu(m)}\prod_jm_j!$: rotations contribute $p^{\nu(m)}$, and reordering
components of height $j$ contributes $m_j!$. These lifts form one $S_m$-orbit:
compare two forests by matching heights, then match leaves through arbitrarily
chosen orders. Quotienting by their labels therefore shows transitivity of
$S_m$ on $\F(m)$ and the asserted stabilizer, proving \eqref{digit}.\end{proof}

\begin{lemma}\label{K} Let $K(a)$ count permutations of $[pa]$ all of whose
cycles have length $p$. For $a\ge1$,
\begin{equation}\label{Kexact}
 K(a)=\binom{pa-1}{p-1}(p-1)!\,K(a-1)=\frac{(pa)!}{p^aa!},\qquad
 K(a)\equiv(-1)^a\pmod p.
\end{equation}\end{lemma}
\begin{proof} To count, choose the other members and the cyclic order of the
cycle containing $1$, then finish the remainder. Conjugate the set by the cycle
$(1\cdots p)$. A fixed permutation preserves that cycle in one of $p-1$ ways
and has arbitrary $p$-cycle structure on the remaining points. The nonfixed
orbits have size $p$, so $K(a)\equiv(p-1)K(a-1)\equiv-K(a-1)$, with $K(0)=1$.\end{proof}

\begin{lemma}\label{recursion} If $a=\lfloor m/p\rfloor$ and $r=m-pa$,
\begin{equation}\label{forest}
 |\F(m)|=\binom mrK(a)|\F(a)|,\qquad
 |\F(m)|\equiv(-1)^{\nu(m)}\pmod p.
\end{equation}\end{lemma}
\begin{proof} This is a direct algebraic consequence of \eqref{digit} and
\eqref{Kexact}, not a suppressed new valuation argument. Base-$p$ digit
identities $\nu(a)=\nu(m)-a$ and $a_j=m_{j+1}$ give
\begin{align*}
\binom mrK(a)|\F(a)|
&=\frac{m!}{r!(pa)!}\frac{(pa)!}{p^a a!}
 \frac{a!}{p^{\nu(a)}\prod a_j!}\\
&=\frac{m!}{r!\,p^{\nu(m)}\prod_{j\ge1}m_j!}=|\F(m)|.
\end{align*}
For the congruence, rotate separately inside $a$ consecutive blocks of $p$
points of $[m]$, holding the last $r$ points. A fixed $r$-subset contains only
entire moving blocks or fixed points. As $r<p$, it is uniquely the fixed points.
All other subset orbits have size $p$, hence $\binom mr\equiv1\pmod p$.
If $m<p$, then $\nu(m)=0$ and $|\F(m)|=1$; for $a=0$ the count $K(0)=1$.
Lemma~\ref{K}, the proven recurrence and $\nu(m)=a+\nu(a)$ complete induction
on $m$.\end{proof}

\section{The theorem}
\begin{theorem}\label{main} For every prime $p$, $k\ge1$ and $n\ge1$,
with $m=(p^k-1)n+1$,
\begin{equation}\label{maineq}
 \boxed{C_n^{(q)}\prod_jm_j!\equiv(-p)^t\pmod{p^{t+1}}}
\end{equation}
in $\mathbb Z$, and $v_p(C_n^{(q)})=t=(s_p(m)-1)/(p-1)$.\end{theorem}
\begin{proof} By Lemmas~\ref{L1} and \ref{stabilizer},
\[
 p^{N(m)}L(m)=p^{\nu(m)}\prod_jm_j!|\F(m)|C_n^{(q)}.
\]
Cancelling $p^\nu$ and recalling $t=N-\nu$ therefore yields the integer
identity
\[
 C_n^{(q)}\prod_jm_j!=p^t\frac{L(m)}{|\F(m)|}.
\]
Both $L(m)$ and $|\F(m)|$ are $p$-units by Lemmas~\ref{L2} and \ref{recursion}.
Respectively they are congruent to $(-1)^N$ and $(-1)^\nu$ modulo $p$.
Since
\[
 L(m)\equiv(-1)^{N-\nu}|\F(m)|\pmod p,
\]
their $p$-adic unit quotient is $(-1)^t$ modulo $p$, so its integer expression
has the form $(-1)^t+p\lambda$. Multiplication by $p^t$ proves \eqref{maineq}.
Each $m_j<p$ makes $\prod m_j!$ a $p$-unit, so $v_p(C_n^{(q)})=t$.\end{proof}
\begin{remark}\label{qualification} Here $v_p(m!)$ \emph{is} computed along the
way but from complete-tree symmetry rather than from Legendre's formula. Thus
this note does not meet an absolute prohibition on factorial valuations;
whether that satisfies the intended interpretation is for humans to judge.
In the original skeleton proof, these digit trees are atoms, with the remaining
$t$ internal nodes forming the carry skeleton. A carry replaces $p$ atoms of
one height by one atom of the next height, read modulo $k$; label cycling
and collapsing translate that carry into Lemma~\ref{cherry}.\end{remark}

\section{Sand on a ring and on a line}
Take $a_c\ge0$ grains at sites $c=0,1,\ldots$ with finite support, or $c\in
\mathbb Z/k$. A site with $a_c\ge p$ topples to $(a_c-p,a_{c+1}+1)$, modulo $k$
on the ring. At the one-site ring, grain $1$ returns to the same site.
\begin{proposition}\label{sand} Every legal sequence is finite, and all maximal
legal sequences from a state have the same final state and per-site toppling
counts. From $m$ grains at the origin:
\begin{enumerate}
\item On the half-line the terminal state is the base-$p$ digits of $m$;
$x_c=\lfloor m/p^{c+1}\rfloor$ grains topple at site $c$, totalling $\nu(m)$.
\item On the ring the conserved quantity is $w(a)=\sum a_cp^c\pmod{q-1}$. The
terminal state is that of the unique $r\in\{1,\ldots,q-1\}$ congruent to $m$
modulo $q-1$ (with its $k$ digits zero-extended). For $q=2$ the final grain is
always $1$. For $m=(q-1)n+1$, the terminal state is $e_0$ and topplings total
$N(m)$.
\item Starting from the line's terminal digits folded into $k$ sites, exactly
$t$ ring topplings remain. Therefore $v_p(C_n^{(q)})$ is the difference of the
topplings on the ring and on the line, both starting with $m$ at site zero.
\item If $S(a)_c=a_{c+1}$ on the ring, then $\mathbb Z^k/\operatorname{im}(pI-S)
\cong\mathbb Z/(q-1)\mathbb Z$.
\end{enumerate}\end{proposition}
\begin{proof} Firing lowers the total grain sum by $p-1$, so every legal
sequence is finite. For abelianness, induct upon the sum and consider two
maximal legal words. Their first firings can be interchanged: legal firings
on all one-site initial states commute straightforwardly on the generic line,
and on a ring immediate wrap firings also commute because adding grain never
makes a legal firing illegal. If the first firing of the first word appears
somewhere in the second, move it to the front using that local diamond
property. If it never appears, it remains legal after the whole second word,
contradicting maximality. Cancel that common initial firing and induct.

On the half-line, grain preservation gives $m-p x_0=b_0$, and for $c\ge1$,
$x_{c-1}-p x_c=b_c$ where $b_c$ is the terminal amount. Hence
$M_0=m$ and $M_c=x_{c-1}$ gives $M_c=p x_c+b_c$. If
$M_c=\lfloor m/p^c\rfloor$, then $0\le b_c<p$ forces $b_c=m_c$ and
$x_c=M_{c+1}=\lfloor m/p^{c+1}\rfloor$. This proves (1).

On the ring, $p^k\equiv1\pmod{q-1}$ conserves $w(a)$. The terminal digits
are $0,\ldots,p-1$ and their integer weight lies between $1$ and $q-1$ for
positive input, namely the stated $r$. It remains to locate the remainder:
it is a cyclic translate of $b=(b_0,\ldots,b_{k-1})$ on $k$ nodes (with all
higher digits zero) if and only if there are nonnegative $x$ and positive
quotient $\ell$ with
\[
 w(b)-r=\sum_j(b_j-r_j)p^j=\ell(q-1).
\]
The difference of two stable vectors has $|w|\le q-2$, since neither has
all $p-1$ digits: $q-1$ itself is not a terminal remainder for positive
input reachable from arbitrary positive grain sum? Actually that stable state
can be reached from inputs, and weight congruence allows only zero difference
unless the all-$(p-1)$ state is one of them. So the following essential
boundary exception requires care. \GAP{Prove the positional uniqueness of
the terminal base-$p$ remainder, with correct handling of the boundary
$q-1$; the weighted invariant alone is insufficient. I have not verified an
oriented cyclic-shift invariant. The simulation is not a replacement for
this argument.}

For the special initial vector $(m,0,\ldots,0)$, $m=(q-1)n+1$, there is one
state which completely avoids this general positional issue: its candidate
nonnegative toppling vector is
\[
 x_c=(p^{k-c}-1)n,\qquad c=0,\ldots,k-1.
\]
It satisfies $p x_0-x_{k-1}=m-1$ and $x_{c-1}-p x_c=0$ for $c>0$. Therefore
with $A=pI-S$, $Ax=me_0-e_0$; the candidate terminal state is $e_0$.
Because an actual legal sequence exists, it too produces a vector $y$ with
$Ay=me_0-b$. If $b$ is stable, obviously $w(b)=1$ is one of its valid forms,
but $A^{-1}(me_0-e_0)=x$ and all forms $b$ are determined solely by toppling
counts. Since for any transformed stable state $b$ the equation
$me_0-b=Ay$ has a unique rational solution, its difference $x-y$ must yield
$b-e_0$ in the image. The nonnegative character and digits of $b$ can be
considered directly using $A^{-1}(b-e_0)$: if $b=e_c$ for some $c\ne0$, its
integer weight $p^c$ is not congruent to $1$ modulo $q-1$ because
$1<p^c<q$; this finally forces $b=e_0$, and hence also $x=y$. At this point
this argument does not require ring root-shift uniqueness for arbitrary $m$:
stability, sum and congruence at the actual final state, together with the
specific $r=1$ single-digit candidate, do imply that no alternate terminal
position is compatible with weight conservation. The site-count formula
$x_c=(p^{k-c}-1)n$ is also uniquely recovered from the terminal equation.
It sums to $n((q-1)-k)/(p-1)=N(m)$. The total toppling count can moreover be
obtained without identifying the site vector, since total grains begin at $m$
and end at $s_p(m)$: it is exactly $N(m)$.

Starting from the folded line digits uses total $s_p(m)$ grains, so its
remaining ring firings number $(s_p(m)-1)/(p-1)=t$, using the special-case
endpoint just proved; its original count from $m$ is the difference
$N(m)-\nu(m)$.

For (4), the characteristic calculation is
$\det(pI-S)=q-1$ by expanding the cycle permutation term. Deleting row $1$
and column $k$ gives a triangular minor of determinant $\pm1$. Hence all
Smith divisors except the last are $1$, and the last must be $q-1$. The integer
cokernel is therefore $\mathbb Z/(q-1)\mathbb Z$.\end{proof}
\end{document}
